% Auto-generated from book/C_lab_code_reference.md — do not edit by hand
% Regenerated by scripts/md_to_latex.py

\chapter{Laboratory Code Reference}
\label{ch:app_c_labs}

Full lab listings moved out of the main chapter flow. Chapters keep short ``call'' versions and exercises; use this appendix when typing or copying longer scripts.

\textbf{Setup (all labs):}

\begin{lstlisting}[style=qga,language=python]
import sys
from pathlib import Path
# book helpers
sys.path.insert(0, str(Path.home() / "Projects" / "qga"))
# portal (when needed)
# export PYTHONPATH=~/Projects/kingdom/src:~/Projects/flux_hopf_lib/src
\end{lstlisting}

\bigskip
\noindent\rule{\textwidth}{0.4pt}
\bigskip

\section{C.1 Chapters 1--2 (foundations)}
\label{ch:app_c_labs:c-1-chapters-1-2-foundations}

\subsection{Lab 1.A --- Multiplication table}
\label{ch:app_c_labs:lab-1-a-multiplication-table}

\begin{lstlisting}[style=qga,language=python]
from kingdom.core.quaternion import Quaternion

i = Quaternion(0, 1, 0, 0)
j = Quaternion(0, 0, 1, 0)
k = Quaternion(0, 0, 0, 1)
print(i.multiply(j))  # ~ k
print(j.multiply(i))  # ~ -k
print(i.multiply(i))  # ~ -1
\end{lstlisting}

\subsection{Lab 1.B --- Norm multiplicativity}
\label{ch:app_c_labs:lab-1-b-norm-multiplicativity}

\begin{lstlisting}[style=qga,language=python]
import numpy as np
from kingdom.core.quaternion import Quaternion

rng = np.random.default_rng(0)

def rand_q():
    a = rng.normal(size=4)
    return Quaternion(*a)

q1, q2 = rand_q(), rand_q()
prod = q1.multiply(q2)
n1, n2 = q1.norm() ** 2, q2.norm() ** 2
print(abs(prod.norm() ** 2 - n1 * n2))
\end{lstlisting}

\subsection{Lab 1.D --- Double cover}
\label{ch:app_c_labs:lab-1-d-double-cover}

\begin{lstlisting}[style=qga,language=python]
import numpy as np
from kingdom.core.quaternion import Quaternion

q = Quaternion.from_axis_angle(np.array([0.0, 0.0, 1.0]), np.pi / 3)
mq = Quaternion(-q.w, -q.x, -q.y, -q.z)
v = Quaternion(0, 1, 0, 0)
rq = q.multiply(v).multiply(q.inverse())
rm = mq.multiply(v).multiply(mq.inverse())
print(rq, rm)  # same rotation
\end{lstlisting}

\subsection{Lab 2.A--B --- Hopf image and fiber}
\label{ch:app_c_labs:lab-2-a-b-hopf-image-and-fiber}

\begin{lstlisting}[style=qga,language=python]
from kingdom.core.quaternion import Quaternion
from kingdom.core.hopf import hopf_map_quaternion, sample_fiber
import numpy as np

q = Quaternion(0.5, 0.5, 0.5, 0.5).normalize()
print(q.hopf_image(), hopf_map_quaternion(q.w, q.x, q.y, q.z))

fiber = sample_fiber(eta=0.4, xi1=0.3, n_points=64)
print(sorted(fiber.keys()))
X = np.stack([fiber["x1"], fiber["x2"], fiber["x3"], fiber["x4"]], axis=1)
print(np.linalg.norm(X, axis=1)[:5])
\end{lstlisting}

\subsection{Lab 2.E --- Fiber family}
\label{ch:app_c_labs:lab-2-e-fiber-family}

\begin{lstlisting}[style=qga,language=python]
from kingdom.core.hopf import sample_fiber_family
family = sample_fiber_family(n_fibers=12, n_points=160)
print(len(family), family[0]["base_y1"])
\end{lstlisting}

\bigskip
\noindent\rule{\textwidth}{0.4pt}
\bigskip

\section{C.2 Chapters 3--4 (lattice and symmetries)}
\label{ch:app_c_labs:c-2-chapters-3-4-lattice-and-symmetries}

\subsection{Lab 3.A--C --- Hurwitz, adjacency, gauge}
\label{ch:app_c_labs:lab-3-a-c-hurwitz-adjacency-gauge}

\begin{lstlisting}[style=qga,language=python]
from lib.hopf_lattice import (
    HURWITZ_UNITS, hopf_project_points, sample_angle_lattice,
    candidate_adjacency, left_multiply, right_multiply,
)
import numpy as np

print(len(HURWITZ_UNITS), np.allclose(np.sum(HURWITZ_UNITS**2, axis=1), 1.0))
base = hopf_project_points(HURWITZ_UNITS)
print("distinct base ~", len(np.unique(np.round(base, 5), axis=0)))

pts = sample_angle_lattice(n_eta=3, n_xi1=8, n_xi2=8)
along, inter = candidate_adjacency(pts, base_angle_thresh=0.5)
print(len(pts), len(along), len(inter))

i = np.array([0.0, 1.0, 0.0, 0.0])
moved = left_multiply(HURWITZ_UNITS, i)
base0, baseL = hopf_project_points(HURWITZ_UNITS), hopf_project_points(moved)
same_set = np.array_equal(np.sort(np.round(base0, 6), 0), np.sort(np.round(baseL, 6), 0))
print("base set invariant:", same_set)
print("total-space moved:", not np.allclose(moved, HURWITZ_UNITS))
\end{lstlisting}

\subsection{Lab 3.D --- Two-gyro dynamics}
\label{ch:app_c_labs:lab-3-d-two-gyro-dynamics}

\begin{lstlisting}[style=qga,language=python]
from kingdom.core.lattice import LatticeConfig
from kingdom.simulations.lattice import TwoGyroLattice, run_lattice_comparison

stable, chaotic = run_lattice_comparison(frames=80, n_sites=48)
print(stable.stability_score, stable.total_bursts, chaotic.total_bursts)
\end{lstlisting}

\subsection{Lab 4.A--B --- Gauge sequences and orbits}
\label{ch:app_c_labs:lab-4-a-b-gauge-sequences-and-orbits}

\begin{lstlisting}[style=qga,language=python]
from lib.hopf_lattice import (
    HURWITZ_UNITS, permutes_hurwitz_units, left_multiply,
    hopf_project_points, orbit_of_point, phase_unit,
)
import numpy as np

i = np.array([0.0, 1.0, 0.0, 0.0])
print(permutes_hurwitz_units(i, side="L"))
seq = [("R", phase_unit(np.pi / 3)), ("L", i)]
orbit = orbit_of_point(HURWITZ_UNITS[3], seq, max_periods=24, tol=1e-6)
print(len(orbit), np.linalg.norm(orbit[-1] - orbit[0]))
\end{lstlisting}

\subsection{Lab 4.D --- Flux push-forward + OP1 equivariance}
\label{ch:app_c_labs:lab-4-d-flux-push-forward-op1-equivariance}

\begin{lstlisting}[style=qga,language=python]
from lib.hopf_lattice import (
    sample_angle_lattice, candidate_adjacency, discrete_flux_cycle,
    left_multiply, transform_flux, adjacency_equivariance_score,
)
import numpy as np

pts = sample_angle_lattice(2, 6, 8)
along, inter = candidate_adjacency(pts, base_angle_thresh=0.55, fiber_phase_bins=8)
edges = along[:8] if along else inter[:8]
Phi = discrete_flux_cycle(edges, value=1)
i = np.array([0.0, 1.0, 0.0, 0.0])
imap = {k: k for k in range(len(pts))}
print(len(transform_flux(Phi, imap)) // 2)
print(adjacency_equivariance_score(pts, i, side="L", base_angle_thresh=0.55, fiber_phase_bins=8))
\end{lstlisting}

\bigskip
\noindent\rule{\textwidth}{0.4pt}
\bigskip

\section{C.3 Chapters 5--7 (topographs, classification, Z-map)}
\label{ch:app_c_labs:c-3-chapters-5-7-topographs-classification-z-map}

\subsection{Lab 5.A--D --- Topographs}
\label{ch:app_c_labs:lab-5-a-d-topographs}

\begin{lstlisting}[style=qga,language=python]
from lib.hopf_lattice import sample_angle_lattice, candidate_adjacency, phase_unit
from lib.flux_topograph import (
    build_flux_topograph, detect_separators, arithmetic_progression_residuals,
    periodicity_score, separator_equivariance_score,
)
import numpy as np

pts = sample_angle_lattice(4, 12, 12)
along, inter = candidate_adjacency(pts, base_angle_thresh=0.5, fiber_phase_bins=12)
topo = build_flux_topograph(pts, edges=along + inter, functional="hopf_height")
seps = detect_separators(topo, mode="sign")
print(len(seps), sum(len(c) for c in seps), arithmetic_progression_residuals(topo))

i = np.array([0.0, 1.0, 0.0, 0.0])
print(periodicity_score(topo, [("L", i), ("R", phase_unit(np.pi / 2))], max_periods=8))
print(separator_equivariance_score(topo, [("L", i)]))
\end{lstlisting}

\subsection{Lab 6.A--D --- Classification}
\label{ch:app_c_labs:lab-6-a-d-classification}

\begin{lstlisting}[style=qga,language=python]
from lib.flux_topograph import (
    build_flux_topograph, classify_topograph_type, enumerate_reduced,
    class_number_analogue, equivalence_distance, apply_gauge_to_topograph,
)
from lib.hopf_lattice import sample_angle_lattice, candidate_adjacency
import numpy as np

pts = sample_angle_lattice(3, 10, 10)
along, inter = candidate_adjacency(pts, base_angle_thresh=0.5, fiber_phase_bins=10)
topo = build_flux_topograph(pts, edges=along + inter, functional="hopf_height")
print(classify_topograph_type(topo))
print("n reduced", len(enumerate_reduced(topo, dedup_tol=0.05)))
print(class_number_analogue(topo, dedup_tol=0.05)["class_number_analogue"])
i = np.array([0.0, 1.0, 0.0, 0.0])
print(equivalence_distance(topo, apply_gauge_to_topograph(topo, [("L", i)])))
\end{lstlisting}

\subsection{Lab 7.A--E --- Z-map}
\label{ch:app_c_labs:lab-7-a-e-z-map}

\begin{lstlisting}[style=qga,language=python]
from kingdom.core.flux_flywheel import map_z_to_flywheel, map_z_to_flywheel_extended
from lib.flux_topograph import stability_landscape_z

for z in [2, 10, 26, 79, 118]:
    m = map_z_to_flywheel(z)
    print(z, m["stability_score"], m["stability_class"], m.get("is_noble_gas"))

landscape = stability_landscape_z(z_range=(1, 50))
high = [r for r in landscape if (r["stability_score"] or 0) >= 7.5]
print([(r["Z"], r["stability_score"]) for r in high])

for z in [2, 26]:
    ext = map_z_to_flywheel_extended(z)
    print({k: ext.get(k) for k in [
        "stability_score", "is_noble_gas", "alignment_stability_pts",
        "real_ionization_energy_eV", "model_vs_reality_alignment",
    ]})
\end{lstlisting}

\bigskip
\noindent\rule{\textwidth}{0.4pt}
\bigskip

\section{C.4 Chapters 8--10 (composition, algebras, validation)}
\label{ch:app_c_labs:c-4-chapters-8-10-composition-algebras-validation}

\subsection{Lab 8.A--D --- Composition}
\label{ch:app_c_labs:lab-8-a-d-composition}

\begin{lstlisting}[style=qga,language=python]
from lib.hopf_lattice import sample_angle_lattice, candidate_adjacency
from lib.flux_topograph import build_flux_topograph, class_number_analogue
from lib.composition import (
    compose_flywheels, reduce_composition, composition_table,
    class_group_analogue, is_associative_up_to_equivalence,
)

pts = sample_angle_lattice(2, 6, 6)
along, inter = candidate_adjacency(pts, base_angle_thresh=0.55, fiber_phase_bins=6)
edges = along + inter
t1 = build_flux_topograph(pts, edges=edges, functional="hopf_height")
t2 = build_flux_topograph(pts, edges=edges, functional="hopf_y1")
print(reduce_composition(compose_flywheels(t1, t2, method="value_sum"))["classification"]["type"])

cg = class_group_analogue(t1, dedup_tol=0.1, samples=10)
print(cg["order"], cg["structure"], cg["associativity"])
reps = cg["representatives"]
print(composition_table(reps, dedup_tol=0.1)["closure_fraction"])
print(is_associative_up_to_equivalence(reps, samples=20, tol=0.1))
\end{lstlisting}

\subsection{Lab 9.A--E --- Quaternion algebras}
\label{ch:app_c_labs:lab-9-a-e-quaternion-algebras}

\begin{lstlisting}[style=qga,language=python]
from lib.quaternion_algebra import (
    QuaternionAlgebra, HurwitzOrder, LipschitzOrder,
    left_ideal_class_group, hilbert_symbol,
)

A = QuaternionAlgebra(-1, -1)
print(A.presentation(), A.ramified_places(), A.is_definite())
O = HurwitzOrder()
print(O.n_units(), O.is_euclidean(), O.is_maximal())
print(left_ideal_class_group(O).order, left_ideal_class_group(O).method)
print({p: hilbert_symbol(-1, -1, p) for p in [2, 3, 5, "inf"]})
\end{lstlisting}

\subsection{Lab 10.B-style --- Table T4 demo}
\label{ch:app_c_labs:lab-10-b-style-table-t4-demo}

\begin{lstlisting}[style=qga,language=python]
from lib.validation import (
    run_table_t4_demo, combine_p_values_fisher, bonferroni_threshold,
    proximity_to_wg, WG_350_OVER_PI, default_hypotheses, table_t4_checklist,
)

print("W_g", WG_350_OVER_PI)
print(len(table_t4_checklist()), [h.name for h in default_hypotheses()])
demo = run_table_t4_demo(seed=1, alpha=0.01)
print(demo["bonferroni_threshold"], demo["decision"], demo["disclaimer"])
print(proximity_to_wg([111.4, 111.5, 110.0, 111.408]))
\end{lstlisting}

\bigskip
\noindent\rule{\textwidth}{0.4pt}
\bigskip

\section{C.5 Common pitfalls}
\label{ch:app_c_labs:c-5-common-pitfalls}

\begin{center}
\small
\begin{tabular}{@{}lp{0.55\textwidth}@{}}
\toprule
Symptom & Fix \\
\midrule
\texttt{ModuleNotFoundError: lib} & \texttt{sys.path} or run from repo root \\
\texttt{ModuleNotFoundError: kingdom} & Set \texttt{PYTHONPATH} or use kingdom venv \\
\texttt{magic\_flag} KeyError & Use \texttt{stability\_score} / \texttt{is\_noble\_gas} \\
\texttt{step} AttributeError & Use \texttt{step\_frame()} \\
Slow \texttt{enumerate\_reduced} & Smaller lattices (\texttt{n\_eta\textbackslash{}(\textbackslash{}leq\textbackslash{})3}, \texttt{n\_xi\textbackslash{}(\textbackslash{}leq\textbackslash{})10}) \\
\bottomrule
\end{tabular}
\end{center}


\bigskip
\noindent\rule{\textwidth}{0.4pt}
\bigskip


